\documentclass[12pt, a4paper]{article}
\usepackage{../../../myconfig}

\begin{document}

\titlebox{CHUYÊN ĐỀ: OXYZ}{BÀI TẬP VỀ NHÀ TOÀN DIỆN\\ (CÓ DÒNG KẺ TRÌNH BÀY BÀI TẬP)}{Oxyz: Bài tập về nhà (Lines)}

% ==========================================================
% DẠNG 1: TỌA ĐỘ ĐIỂM, VECTƠ, HÌNH CHIẾU & ĐỐI XỨNG (25 CÂU - 4 DÒNG)
% ==========================================================
\noindent\textbf{\textcolor{deepblue}{\large DẠNG 1: TỌA ĐỘ ĐIỂM, VECTƠ, HÌNH CHIẾU \& ĐỐI XỨNG (CÂU 1 -- 25)}}
\par\vspace{0.2cm}

\questionLinesManual{4}{1}{Trong không gian với hệ trục tọa độ $Oxyz$, cho véctơ $\vec{u}=2\vec{i}-3\vec{j}+4\vec{k}$. Tọa độ của véctơ $\vec{u}$ là:}{%
    \begin{tasks}(4)
        \task $(2;3;4)$
        \task $(4;-3;2)$
        \task $(2;-3;4)$
        \task $(-3;2;4)$
    \end{tasks}
}

\questionLinesManual{4}{2}{Trong không gian $Oxyz$, cho vectơ $\vec{u}=2\vec{i}-5\vec{k}$. Tọa độ của vectơ $\vec{u}$ là:}{%
    \begin{tasks}(4)
        \task $(0;2;-5)$
        \task $(2;0;5)$
        \task $(2;-5;0)$
        \task $(2;0;-5)$
    \end{tasks}
}

\questionLinesManual{4}{3}{Trong không gian $Oxyz$, cho $\vec{a}=-\vec{i}+3\vec{j}-5\vec{k}$. Tọa độ của vectơ $\vec{a}$ là:}{%
    \begin{tasks}(4)
        \task $(-1;3;-5)$
        \task $(1;-3;5)$
        \task $(-5;3;-1)$
        \task $(5;-3;1)$
    \end{tasks}
}

\questionLinesManual{4}{4}{Trong không gian $Oxyz$, cho $\vec{OA}=6\vec{j}+4\vec{i}-3\vec{k}$. Tọa độ của điểm $A$ là:}{%
    \begin{tasks}(4)
        \task $(4;6;-3)$
        \task $(-4;-6;3)$
        \task $(-6;-4;3)$
        \task $(6;4;-3)$
    \end{tasks}
}

\questionLinesManual{4}{5}{Trong không gian $Oxyz$, cho điểm $M(-2;5;7)$. Tọa độ của vectơ $\vec{MO}$ là:}{%
    \begin{tasks}(4)
        \task $(-2;-5;7)$
        \task $(-2;5;7)$
        \task $(2;-5;-7)$
        \task $(2;5;7)$
    \end{tasks}
}

\questionLinesManual{4}{6}{Trong không gian $Oxyz$, cho điểm $A(1;2;-1)$. Tọa độ hình chiếu vuông góc của $A$ trên mặt phẳng $(Oyz)$ là:}{%
    \begin{tasks}(4)
        \task $(0;2;-1)$
        \task $(1;0;0)$
        \task $(1;2;0)$
        \task $(1;0;-1)$
    \end{tasks}
}

\questionLinesManual{4}{7}{Trong không gian $Oxyz$, hình chiếu vuông góc của điểm $M(2;-2;1)$ trên mặt phẳng $(Oxy)$ có tọa độ là:}{%
    \begin{tasks}(4)
        \task $(2;-2;0)$
        \task $(2;0;1)$
        \task $(0;-2;1)$
        \task $(0;0;1)$
    \end{tasks}
}

\questionLinesManual{4}{8}{Trong không gian $Oxyz$, cho điểm $A(1;-2;5)$. Hình chiếu vuông góc của điểm $A$ lên trục $Ox$ là:}{%
    \begin{tasks}(4)
        \task $(0;-2;5)$
        \task $(0;0;5)$
        \task $(0;-2;0)$
        \task $(1;0;0)$
    \end{tasks}
}

\questionLinesManual{4}{9}{Trong không gian $Oxyz$, hình chiếu vuông góc của điểm $M(-2;3;4)$ lên trục $Oy$ là điểm nào?}{%
    \begin{tasks}(4)
        \task $M_1(-2;0;0)$
        \task $M_2(0;3;0)$
        \task $M_3(0;0;4)$
        \task $M_4(-2;0;4)$
    \end{tasks}
}

\questionLinesManual{4}{10}{Trong không gian $Oxyz$, cho điểm $A(1;2;-3)$. Hình chiếu vuông góc của $A$ lên mặt phẳng $(Oxy)$ có tọa độ là:}{%
    \begin{tasks}(4)
        \task $(1;0;0)$
        \task $(0;2;-3)$
        \task $(1;0;-3)$
        \task $(1;2;0)$
    \end{tasks}
}

\questionLinesManual{4}{11}{Hai điểm $M$ và $M'$ phân biệt đối xứng nhau qua mặt phẳng $(Oxy)$. Phát biểu nào sau đây là đúng?}{%
    \begin{tasks}(1)
        \task Hai điểm $M$ và $M'$ có cùng tung độ và cao độ.
        \task Hai điểm $M$ và $M'$ có cùng hoành độ và cao độ.
        \task Hai điểm $M$ và $M'$ có hoành độ đối nhau.
        \task Hai điểm $M$ và $M'$ có cùng hoành độ và tung độ.
    \end{tasks}
}

\questionLinesManual{4}{12}{Trong không gian $Oxyz$, cho điểm $A(-5;2;3)$ và $B$ là điểm đối xứng với $A$ qua trục $Oy$. Độ dài đoạn thẳng $AB$ bằng:}{%
    \begin{tasks}(4)
        \task $\sqrt{34}$
        \task $2\sqrt{38}$
        \task $2\sqrt{34}$
        \task $\sqrt{38}$
    \end{tasks}
}

\questionLinesManual{4}{13}{Trong không gian $Oxyz$, cho hai điểm $A(-2;3;5)$ và $B$ là điểm đối xứng với $A$ qua trục $Oz$. Độ dài đoạn thẳng $AB$ bằng:}{%
    \begin{tasks}(4)
        \task $2\sqrt{34}$
        \task $\sqrt{13}$
        \task $\sqrt{34}$
        \task $2\sqrt{13}$
    \end{tasks}
}

\questionLinesManual{4}{14}{Trong không gian $Oxyz$, cho điểm $M(1;3;-2)$. Gọi $M'$ là hình chiếu vuông góc của $M$ trên trục $Oz$. Tọa độ của vectơ $\vec{MM'}$ là:}{%
    \begin{tasks}(4)
        \task $(1;3;0)$
        \task $(0;0;2)$
        \task $(-1;-3;0)$
        \task $(0;0;-2)$
    \end{tasks}
}

\questionLinesManual{4}{15}{Trong không gian $Oxyz$, cho hai điểm $A(1;1;-2)$ và $B(2;2;1)$. Vectơ $\vec{AB}$ có tọa độ là:}{%
    \begin{tasks}(4)
        \task $(1;1;3)$
        \task $(3;1;1)$
        \task $(-1;-1;-3)$
        \task $(3;3;-1)$
    \end{tasks}
}

\questionLinesManual{4}{16}{Trong không gian $Oxyz$, cho hai điểm $A(-1;2;1)$ và $B(2;1;-3)$. Tọa độ của vectơ $\vec{AB}$ là:}{%
    \begin{tasks}(4)
        \task $\left(\dfrac{1}{2}; \dfrac{3}{2}; -1\right)$
        \task $(-3;1;4)$
        \task $(3;-1;-4)$
        \task $(1;3;-2)$
    \end{tasks}
}

\questionLinesManual{4}{17}{Trong không gian $Oxyz$, cho hai điểm $A(-2;1;0)$, $B(3;-2;1)$. Tọa độ của vectơ $\vec{AB}$ là:}{%
    \begin{tasks}(4)
        \task $(-5;3;-1)$
        \task $(5;-3;1)$
        \task $(1;-1;1)$
        \task $(-1;1;-1)$
    \end{tasks}
}

\questionLinesManual{4}{18}{Trong không gian $Oxyz$, cho hai điểm $A(7;-4;-2)$ và $B(-9;-9;7)$. Tọa độ của vectơ $\vec{AB}$ là:}{%
    \begin{tasks}(4)
        \task $(-16;-5;9)$
        \task $(-2;-13;5)$
        \task $(16;5;-9)$
        \task $\left(-1; -\dfrac{13}{2}; \dfrac{5}{2}\right)$
    \end{tasks}
}

\questionLinesManual{4}{19}{Trong không gian $Oxyz$, cho hai điểm $A(2;-1;-1)$, $B(-3;2;-2)$. Tọa độ của vectơ $\vec{AB}$ là:}{%
    \begin{tasks}(4)
        \task $(5;-3;-1)$
        \task $(5;-3;1)$
        \task $(-5;1;-1)$
        \task $(-5;3;-1)$
    \end{tasks}
}

\questionLinesManual{4}{20}{Trong không gian $Oxyz$, cho các vectơ $\vec{OA}=-\vec{i}+\vec{k}$ và $\vec{OB}=\vec{i}-\vec{j}+2\vec{k}$. Tọa độ vectơ $\vec{AB}$ là:}{%
    \begin{tasks}(4)
        \task $(-2;1;-1)$
        \task $(2;-1;1)$
        \task $(0;-1;3)$
        \task $(0;-1;-1)$
    \end{tasks}
}

\questionLinesManual{4}{21}{Trong không gian $Oxyz$, cho $A(1;0;1)$. Tìm tọa độ điểm $C$ thỏa mãn $\vec{AC}=(0;6;1)$:}{%
    \begin{tasks}(4)
        \task $C(-1;6;-1)$
        \task $C(1;6;2)$
        \task $C(1;6;0)$
        \task $C(-1;-6;-2)$
    \end{tasks}
}

\questionLinesManual{4}{22}{Trong không gian $Oxyz$, cho điểm $M(1;2;-3)$ và vectơ $\vec{v}=(2;-1;-2)$ thỏa mãn $\vec{v}=\vec{MN}$. Tọa độ điểm $N$ là:}{%
    \begin{tasks}(4)
        \task $(-1;3;-1)$
        \task $(3;1;-5)$
        \task $(1;-3;1)$
        \task $(-3;-1;5)$
    \end{tasks}
}

\questionLinesManual{4}{23}{Trong không gian $Oxyz$, cho điểm $M(3;4;5)$. Điểm $N$ thỏa mãn $\vec{MN}=-6\vec{i}$. Tọa độ của điểm $N$ là:}{%
    \begin{tasks}(4)
        \task $N(3;-4;-5)$
        \task $N(-3;-4;-5)$
        \task $N(3;4;-5)$
        \task $N(-3;4;5)$
    \end{tasks}
}

\questionLinesManual{4}{24}{Trong không gian $Oxyz$, cho hình lập phương $ABCD.A'B'C'D'$ có cạnh bằng 1, $A$ trùng với gốc tọa độ $O$, các tia $AB, AD, AA'$ lần lượt trùng với $Ox, Oy, Oz$. Tọa độ vectơ $\vec{A'C}$ là:}{%
    \begin{tasks}(4)
        \task $(1;1;-1)$
        \task $(1;1;1)$
        \task $(1;-1;1)$
        \task $(-1;1;1)$
    \end{tasks}
}

\questionLinesManual{4}{25}{Trong không gian $Oxyz$, cho điểm $M(1;-\sqrt{2};\sqrt{3})$. Tìm điểm $M' \in Ox$ sao cho độ dài đoạn thẳng $MM'$ ngắn nhất:}{%
    \begin{tasks}(4)
        \task $M'(-1;0;0)$
        \task $M'(1;0;0)$
        \task $M'(1;0;\sqrt{3})$
        \task $M'(1;-\sqrt{2};0)$
    \end{tasks}
}

\newpage
% ==========================================================
% DẠNG 2: CÁC PHÉP TOÁN VECTƠ, TRUNG ĐIỂM, TRỌNG TÂM & HÌNH BÌNH HÀNH (25 CÂU - 4 DÒNG)
% ==========================================================
\noindent\textbf{\textcolor{deepblue}{\large DẠNG 2: CÁC PHÉP TOÁN VECTƠ, TRUNG ĐIỂM, TRỌNG TÂM \& HÌNH BÌNH HÀNH (CÂU 26 -- 50)}}
\par\vspace{0.2cm}

\questionLinesManual{4}{26}{Trong không gian với hệ tọa độ $Oxyz$, cho $\vec{a}=(1;3;-1)$ và $\vec{b}=(2;3;6)$. Tọa độ của véctơ $\vec{a}+2\vec{b}$ bằng:}{%
    \begin{tasks}(4)
        \task $(5;-9;-11)$
        \task $(5;-9;11)$
        \task $(5;9;11)$
        \task $(-5;9;11)$
    \end{tasks}
}

\questionLinesManual{4}{27}{Cho ba véc tơ $\vec{a}(1;2;3)$, $\vec{b}(2;2;-1)$, $\vec{c}(4;0;-4)$. Toạ độ của véc tơ $\vec{d}=\vec{a}-\vec{b}+2\vec{c}$ là:}{%
    \begin{tasks}(4)
        \task $\vec{d}(7;0;-4)$
        \task $\vec{d}(-7;0;-4)$
        \task $\vec{d}(7;0;4)$
        \task $\vec{d}(-7;0;4)$
    \end{tasks}
}

\questionLinesManual{4}{28}{Cho $\vec{a}=(1;1;-2)$, $\vec{b}=(-2;1;4)$. Tọa độ của vectơ $\vec{u}=\vec{a}-2\vec{b}$ là:}{%
    \begin{tasks}(4)
        \task $(5;-1;-10)$
        \task $(5;-1;10)$
        \task $(5;1;10)$
        \task $(5;1;-10)$
    \end{tasks}
}

\questionLinesManual{4}{29}{Cho hai vectơ $\vec{u}=(1;-4;0)$ và $\vec{v}=(-1;-2;1)$. Vectơ $\vec{u}+3\vec{v}$ có tọa độ là:}{%
    \begin{tasks}(4)
        \task $(-2;-10;3)$
        \task $(-4;-8;4)$
        \task $(-2;-6;3)$
        \task $(-2;-10;-3)$
    \end{tasks}
}

\questionLinesManual{4}{30}{Cho $\vec{a}=(2;-1;0)$, $\vec{b}=(-1;-3;2)$, $\vec{c}=(-2;-4;-3)$. Tọa độ của $\vec{u}=2\vec{a}-3\vec{b}+\vec{c}$ là:}{%
    \begin{tasks}(4)
        \task $(-3;-7;-9)$
        \task $(-5;-3;9)$
        \task $(5;3;-9)$
        \task $(3;7;9)$
    \end{tasks}
}

\questionLinesManual{4}{31}{Cho ba véctơ $\vec{a}=(5;7;2)$, $\vec{b}=(3;0;4)$, $\vec{c}=(-6;1;-1)$. Tìm tọa độ của véctơ $\vec{m}=3\vec{a}-2\vec{b}+\vec{c}$:}{%
    \begin{tasks}(4)
        \task $\vec{m}=(3;-22;3)$
        \task $\vec{m}=(3;22;3)$
        \task $\vec{m}=(-3;22;-3)$
        \task $\vec{m}=(3;22;-3)$
    \end{tasks}
}

\questionLinesManual{4}{32}{Cho $\vec{a}=(-2;3;1)$, $\vec{b}=(1;-3;4)$. Tìm tọa độ của véctơ $\vec{x}=\vec{b}-\vec{a}$:}{%
    \begin{tasks}(4)
        \task $\vec{x}=(3;-6;3)$
        \task $\vec{x}=(-3;6;-3)$
        \task $\vec{x}=(-1;0;5)$
        \task $\vec{x}=(1;-2;1)$
    \end{tasks}
}

\questionLinesManual{4}{33}{Cho các vectơ $\vec{a}=(1;-1;2)$, $\vec{b}=(2;1;-3)$, $\vec{c}=(0;3;-2)$. Điểm $M(x;y;z)$ thỏa mãn $\vec{OM}+\vec{a}=2\vec{b}-\vec{c}$, tổng $x+y+z$ bằng:}{%
    \begin{tasks}(4)
        \task $3$
        \task $-3$
        \task $4$
        \task $-2$
    \end{tasks}
}

\questionLinesManual{4}{34}{Trong không gian $Oxyz$, cho hai điểm $A(1;0;2)$ và $B(-1;2;0)$. Trung điểm của đoạn thẳng $AB$ có tọa độ là:}{%
    \begin{tasks}(4)
        \task $(0;2;2)$
        \task $(-1;1;-1)$
        \task $(1;1;1)$
        \task $(0;1;1)$
    \end{tasks}
}

\questionLinesManual{4}{35}{Cho hai điểm $A(1;2;3)$, $B(3;8;5)$. Tọa độ trung điểm $I$ của đoạn thẳng $AB$ là:}{%
    \begin{tasks}(4)
        \task $I(2;6;2)$
        \task $I(1;3;1)$
        \task $I(4;10;8)$
        \task $I(2;5;4)$
    \end{tasks}
}

\questionLinesManual{4}{36}{Cho hai điểm $A(1;2;3)$, $B(-2;4;0)$. Trung điểm của đoạn thẳng $AB$ có tung độ bằng:}{%
    \begin{tasks}(4)
        \task $1$
        \task $3$
        \task $2$
        \task $-1$
    \end{tasks}
}

\questionLinesManual{4}{37}{Cho $A(1;1;0)$, $B(3;-1;2)$. Tọa độ điểm $C$ sao cho $B$ là trung điểm của đoạn thẳng $AC$ là:}{%
    \begin{tasks}(4)
        \task $C(4;-3;5)$
        \task $C(-1;3;-2)$
        \task $C(2;0;1)$
        \task $C(5;-3;4)$
    \end{tasks}
}

\questionLinesManual{4}{38}{Cho hai điểm $A(0;-2;1)$, $B(2;-4;3)$. Tìm toạ độ điểm $C$ sao cho $A$ là trung điểm của đoạn thẳng $BC$:}{%
    \begin{tasks}(4)
        \task $C(1;-3;2)$
        \task $C(4;-6;5)$
        \task $C(-2;0;-1)$
        \task $C(2;-2;2)$
    \end{tasks}
}

\questionLinesManual{4}{39}{Cho tam giác $ABC$ với $A(1;3;4)$, $B(2;-1;0)$ và $C(3;1;2)$. Toạ độ trọng tâm $G$ của tam giác $ABC$ là:}{%
    \begin{tasks}(4)
        \task $G\left(3; \dfrac{2}{3}; 3\right)$
        \task $G(2;-1;2)$
        \task $G(2;1;2)$
        \task $G(6;3;6)$
    \end{tasks}
}

\questionLinesManual{4}{40}{Cho tam giác $ABC$ có $A(2;-3;1)$, $B(1;3;-4)$ và $C(3;-3;6)$. Trọng tâm của tam giác $ABC$ có tọa độ là:}{%
    \begin{tasks}(4)
        \task $(2;-1;1)$
        \task $(-6;3;-3)$
        \task $(6;-3;3)$
        \task $(-2;1;-1)$
    \end{tasks}
}

\questionLinesManual{4}{41}{Cho tam giác $ABC$ biết $A(5;-2;0)$, $B(-2;3;0)$, $C(0;2;3)$. Tọa độ trọng tâm $G$ của tam giác $ABC$ là:}{%
    \begin{tasks}(4)
        \task $(1;1;1)$
        \task $(2;0;-1)$
        \task $(1;1;-2)$
        \task $(1;2;1)$
    \end{tasks}
}

\questionLinesManual{4}{42}{Cho hai điểm $A(1;2;-1)$ và $B(2;4;1)$. Trọng tâm của tam giác $OAB$ có tọa độ là:}{%
    \begin{tasks}(4)
        \task $(-1;-2;0)$
        \task $(1;3;0)$
        \task $(1;2;0)$
        \task $(3;6;0)$
    \end{tasks}
}

\questionLinesManual{4}{43}{Cho hai điểm $A(3;4;2)$, $B(-1;-2;2)$ và $G(1;1;3)$ là trọng tâm tam giác $ABC$. Tọa độ điểm $C$ là:}{%
    \begin{tasks}(4)
        \task $C(1;1;5)$
        \task $C(1;3;2)$
        \task $C(0;1;2)$
        \task $C(0;0;2)$
    \end{tasks}
}

\questionLinesManual{4}{44}{Cho bốn điểm $A(1;0;2)$, $B(-2;1;3)$, $C(3;2;4)$, $D(6;9;-5)$. Tọa độ trọng tâm của tứ diện $ABCD$ là:}{%
    \begin{tasks}(4)
        \task $(2;3;1)$
        \task $(-2;3;1)$
        \task $(2;3;1)$
        \task $(2;-3;1)$
    \end{tasks}
}

\questionLinesManual{4}{45}{Cho tam giác $ABC$ có $A(1;1;1)$ và $\vec{AB}+\vec{AC}=(4;0;-6)$. Tọa độ trung điểm $M$ của đoạn thẳng $BC$ là:}{%
    \begin{tasks}(4)
        \task $(5;1;-5)$
        \task $(3;-1;-7)$
        \task $(1;-1;-4)$
        \task $(3;1;-2)$
    \end{tasks}
}

\questionLinesManual{4}{46}{Cho hình bình hành $ABCD$ có $A(-3;1;2)$, $B(-2;4;-1)$, $C(1;-3;3)$. Tọa độ điểm $D$ là:}{%
    \begin{tasks}(4)
        \task $(0;-6;6)$
        \task $(-6;7;-2)$
        \task $(2;0;0)$
        \task $(-4;2;5)$
    \end{tasks}
}

\questionLinesManual{4}{47}{Cho ba điểm $A(1;2;-1)$, $B(2;-1;3)$, $C(-3;5;1)$. Tọa độ điểm $D$ sao cho tứ giác $ABCD$ là hình bình hành là:}{%
    \begin{tasks}(4)
        \task $(-2;8;-3)$
        \task $(-4;8;-3)$
        \task $(-4;8;-5)$
        \task $(-2;2;5)$
    \end{tasks}
}

\questionLinesManual{4}{48}{Cho $A(1;2;-1)$, $B(2;-1;3)$, $C(-2;3;3)$. Điểm $M(a;b;c)$ là đỉnh thứ tư của hình bình hành $ABCM$. Giá trị của $P=a^2+b^2-c^2$ bằng:}{%
    \begin{tasks}(4)
        \task $44$
        \task $6$
        \task $-3$
        \task $-1$
    \end{tasks}
}

\questionLinesManual{4}{49}{Cho các điểm $A(0;-2;-1)$ và $B(1;-1;2)$. Tọa độ điểm $M$ thuộc đoạn thẳng $AB$ sao cho $MA=2MB$ là:}{%
    \begin{tasks}(4)
        \task $\left(\dfrac{2}{3}; -\dfrac{4}{3}; 1\right)$
        \task $\left(\dfrac{1}{2}; -\dfrac{3}{2}; \dfrac{1}{2}\right)$
        \task $(2;0;5)$
        \task $(-1;-3;-4)$
    \end{tasks}
}

\questionLinesManual{4}{50}{Cho hai điểm $B(1;2;-3)$ và $C(7;5;3)$. Tìm tọa độ điểm $E$ thỏa mãn đẳng thức $\vec{CE}=2\vec{EB}$:}{%
    \begin{tasks}(4)
        \task $E(3;9;-1)$
        \task $E(3;3;-1)$
        \task $E\left(3; \dfrac{8}{3}; -\dfrac{8}{3}\right)$
        \task $E(2;3;1)$
    \end{tasks}
}

\newpage
% ==========================================================
% DẠNG 3: TÍCH VÔ HƯỚNG, ĐỘ DÀI, GÓC & BÀI TOÁN THỰC TẾ (25 CÂU - 6 DÒNG)
% ==========================================================
\noindent\textbf{\textcolor{deepblue}{\large DẠNG 3: TÍCH VÔ HƯỚNG, ĐỘ DÀI, GÓC \& BÀI TOÁN THỰC TẾ (CÂU 51 -- 75)}}
\par\vspace{0.2cm}

\questionLinesManual{6}{51}{Trong không gian $Oxyz$, cho các điểm $A(3;0;0)$, $B(0;-4;0)$. Độ dài đoạn thẳng $AB$ bằng:}{%
    \begin{tasks}(4)
        \task $7$
        \task $4$
        \task $3$
        \task $5$
    \end{tasks}
}

\questionLinesManual{6}{52}{Cho hai điểm $A(2;1;3)$, $B(1;-1;5)$. Độ dài đoạn thẳng $AB$ là:}{%
    \begin{tasks}(4)
        \task $5$
        \task $4$
        \task $3$
        \task $6$
    \end{tasks}
}

\questionLinesManual{6}{53}{Cho hai điểm $A(4;2;-2)$ và $B(1;0;1)$. Độ dài đoạn thẳng $AB$ là:}{%
    \begin{tasks}(4)
        \task $22$
        \task $2$
        \task $\sqrt{22}$
        \task $4$
    \end{tasks}
}

\questionLinesManual{6}{54}{Cho hai điểm $A(1;1;-3)$, $B(3;-1;1)$. Gọi $M$ là trung điểm đoạn thẳng $AB$, độ dài đoạn thẳng $OM$ bằng:}{%
    \begin{tasks}(4)
        \task $2\sqrt{6}$
        \task $\sqrt{6}$
        \task $2\sqrt{5}$
        \task $\sqrt{5}$
    \end{tasks}
}

\questionLinesManual{6}{55}{Cho ba điểm $M(1;0;2)$, $N(2;-3;0)$ và $P(2;3;2)$. Chu vi của tam giác $MNP$ bằng:}{%
    \begin{tasks}(4)
        \task $\sqrt{10}+\sqrt{14}+\sqrt{40}$
        \task $\sqrt{5}+\sqrt{13}+\sqrt{17}$
        \task $64$
        \task $\sqrt{20}+\sqrt{22}+\sqrt{24}$
    \end{tasks}
}

\questionLinesManual{6}{56}{Cho $\vec{a}=(1;-2;2)$, $\vec{b}=(-1;2;1)$. Giá trị của tích vô hướng $\vec{a}\cdot\vec{b}$ bằng:}{%
    \begin{tasks}(4)
        \task $3$
        \task $-3$
        \task $2$
        \task $-2$
    \end{tasks}
}

\questionLinesManual{6}{57}{Cho hai vectơ $\vec{a}=(1;2;1)$ và $\vec{b}=(2;-4;-2)$. Khi đó $\vec{a}\cdot\vec{b}$ bằng:}{%
    \begin{tasks}(4)
        \task $12$
        \task $8$
        \task $-12$
        \task $-8$
    \end{tasks}
}

\questionLinesManual{6}{58}{Cho $\vec{a}=(-1;0;2)$ và $\vec{b}=(2;3;-2)$. Giá trị của $\vec{a}\cdot\vec{b}$ bằng:}{%
    \begin{tasks}(4)
        \task $2$
        \task $-4$
        \task $-6$
        \task $-3$
    \end{tasks}
}

\questionLinesManual{6}{59}{Cho hai vectơ $\vec{u}=\vec{i}+3\vec{j}-2\vec{k}$ và $\vec{v}=(2;-1;1)$. Tích vô hướng $\vec{u}\cdot\vec{v}$ bằng:}{%
    \begin{tasks}(4)
        \task $-12$
        \task $5\sqrt{2}$
        \task $-3$
        \task $2\sqrt{21}$
    \end{tasks}
}

\questionLinesManual{6}{60}{Cho $\vec{a}=(2;3;3)$, $\vec{b}=(3;2;-1)$. Khi đó tích vô hướng $\vec{a}\cdot\vec{b}$ bằng:}{%
    \begin{tasks}(4)
        \task $9$
        \task $7$
        \task $3$
        \task $15$
    \end{tasks}
}

\questionLinesManual{6}{61}{Cho hai vectơ $\vec{u}=(2;0;-2)$, $\vec{v}=(-1;-1;6)$. Tích vô hướng $\vec{u}\cdot\vec{v}$ bằng:}{%
    \begin{tasks}(4)
        \task $-14$
        \task $1$
        \task $0$
        \task $4$
    \end{tasks}
}

\questionLinesManual{6}{62}{Cho ba điểm $A(1;-1;2)$, $B(2;0;1)$ và $C(0;-1;3)$. Giá trị của tích vô hướng $\vec{AB}\cdot\vec{AC}$ bằng:}{%
    \begin{tasks}(4)
        \task $0$
        \task $-4$
        \task $-2$
        \task $20$
    \end{tasks}
}

\questionLinesManual{6}{63}{Cho $A(1;3;2)$, $B(1;0;1)$, $C(5;-3;2)$. Biết rằng $\vec{AB}\cdot\vec{AC}=2m$. Giá trị của $m$ là:}{%
    \begin{tasks}(4)
        \task $m=-9$
        \task $m=9$
        \task $m=18$
        \task $m=-18$
    \end{tasks}
}

\questionLinesManual{6}{64}{Côsin của góc giữa hai vectơ $\vec{u}=(10;10;20)$ và $\vec{v}=(10;-20;10)$ là:}{%
    \begin{tasks}(4)
        \task $\dfrac{1}{6}$
        \task $\dfrac{1}{2}$
        \task $-\dfrac{1}{2}$
        \task $-\dfrac{1}{6}$
    \end{tasks}
}

\questionLinesManual{6}{65}{Cho hai vectơ $\vec{a}=(1;-2;1)$, $\vec{b}=(-2;1;1)$. Góc giữa hai vectơ $\vec{a}$ và $\vec{b}$ bằng:}{%
    \begin{tasks}(4)
        \task $60^\circ$
        \task $120^\circ$
        \task $30^\circ$
        \task $-30^\circ$
    \end{tasks}
}

\questionLinesManual{6}{66}{Cho hai vectơ $\vec{m}=(1;-1;1)$ và $\vec{n}=(-1;1;-1)$. Cosin của góc giữa $\vec{m}$ và $\vec{n}$ bằng:}{%
    \begin{tasks}(4)
        \task $1$
        \task $\dfrac{\sqrt{3}}{3}$
        \task $\dfrac{1}{3}$
        \task $-1$
    \end{tasks}
}

\questionLinesManual{6}{67}{Cho $\vec{a}=(1;2;3)$, $\vec{b}=(2;-2;-1)$. Côsin của góc giữa hai vectơ $\vec{a}$ và $\vec{b}$ bằng:}{%
    \begin{tasks}(4)
        \task $-\dfrac{5\sqrt{14}}{42}$
        \task $\dfrac{5\sqrt{14}}{42}$
        \task $\dfrac{5\sqrt{14}}{3}$
        \task $-\dfrac{5\sqrt{14}}{3}$
    \end{tasks}
}

\questionLinesManual{6}{68}{Cho hai véctơ $\vec{a}=(2;1;0)$, $\vec{b}=(-1;0;-2)$. Tính $\cos(\vec{a},\vec{b})$:}{%
    \begin{tasks}(4)
        \task $\dfrac{2}{25}$
        \task $-\dfrac{2}{5}$
        \task $-\dfrac{2}{25}$
        \task $\dfrac{2}{5}$
    \end{tasks}
}

\questionLinesManual{6}{69}{Cho $M(2;3;-1)$, $N(-1;1;1)$, $P(1;m-1;3)$. Tam giác $MNP$ vuông tại $N$ khi $m$ nhận giá trị là:}{%
    \begin{tasks}(4)
        \task $m=3$
        \task $m=0$
        \task $m=1$
        \task $m=2$
    \end{tasks}
}

\questionLinesManual{6}{70}{Cho $A(7;2;8)$ và $M(-4;3;3)$. Tìm tọa độ điểm $B$ biết $\Delta ABC$ vuông tại $C$ và $M$ là tâm đường tròn ngoại tiếp $\Delta ABC$:}{%
    \begin{tasks}(4)
        \task $B(11;-1;5)$
        \task $B(-11;1;-5)$
        \task $B(1;-8;-14)$
        \task $B(-15;4;-2)$
    \end{tasks}
}

\questionLinesManual{6}{71}{Cho hai véctơ $\vec{a}$ và $\vec{b}$ tạo với nhau một góc $120^\circ$, biết $|\vec{a}|=2$ và $|\vec{b}|=4$. Tính $|\vec{a}+\vec{b}|$:}{%
    \begin{tasks}(4)
        \task $\sqrt{8\sqrt{3}+20}$
        \task $2\sqrt{7}$
        \task $2\sqrt{3}$
        \task $6$
    \end{tasks}
}

\questionLinesManual{6}{72}{Cho hai điểm $A(1;1;2)$, $B(-1;3;-9)$. Tọa độ điểm $M$ thuộc trục $Oy$ sao cho tam giác $ABM$ vuông tại $M$ là:}{%
    \begin{tasks}(2)
        \task $M(0; 2\pm 2\sqrt{5}; 0)$
        \task $M(0; 2\pm \sqrt{5}; 0)$
        \task $M(0; 1\pm \sqrt{5}; 0)$
        \task $M(0; 1\pm 2\sqrt{5}; 0)$
    \end{tasks}
}

\questionLinesManual{6}{73}{(Cân bằng lực) Một vật nặng chịu tác dụng của ba lực $\vec{F}_1 = (10; 20; 30)\text{ N}$, $\vec{F}_2 = (-5; 10; -15)\text{ N}$ và $\vec{F}_3$. Biết hệ ở trạng thái cân bằng $\vec{F}_1 + \vec{F}_2 + \vec{F}_3 = \vec{0}$, độ lớn của lực $\vec{F}_3$ xấp xỉ bằng:}{%
    \begin{tasks}(4)
        \task $33{,}9\text{ N}$
        \task $45{,}2\text{ N}$
        \task $25{,}0\text{ N}$
        \task $50{,}0\text{ N}$
    \end{tasks}
}

\questionLinesManual{6}{74}{(Radar kiểm soát không lưu) Một máy bay bay thẳng từ điểm $A(2; 1; 1)\text{ km}$ đến điểm $B(6; 4; 13)\text{ km}$ so với trạm radar tại gốc $O$. Quãng đường bay thẳng $AB$ là:}{%
    \begin{tasks}(4)
        \task $13\text{ km}$
        \task $15\text{ km}$
        \task $12\text{ km}$
        \task $17\text{ km}$
    \end{tasks}
}

\questionLinesManual{6}{75}{(Thiết kế phòng học) Một phòng học dạng hình hộp chữ nhật dài $8\text{ m}$, rộng $6\text{ m}$, cao $3\text{ m}$. Một chiếc đèn được treo tại chính giữa trần nhà. Chọn hệ trục $Oxyz$ có gốc $O$ tại một góc phòng trên mặt sàn $(Oxy)$. Toạ độ điểm treo đèn là:}{%
    \begin{tasks}(4)
        \task $(4;3;3)$
        \task $(4;3;4)$
        \task $(3;4;3)$
        \task $(4;5;3)$
    \end{tasks}
}

\newpage
% ==========================================================
% DẠNG 4: TÍCH CÓ HƯỚNG, DIỆN TÍCH, THỂ TÍCH & BÀI TOÁN PHÂN LOẠI (25 CÂU - 8 DÒNG)
% ==========================================================
\noindent\textbf{\textcolor{deepblue}{\large DẠNG 4: TÍCH CÓ HƯỚNG, DIỆN TÍCH, THỂ TÍCH \& BÀI TOÁN PHÂN LOẠI (CÂU 76 -- 100)}}
\par\vspace{0.2cm}

\questionLinesManual{8}{76}{Trong không gian $Oxyz$, cho $\vec{a}=(1;1;-2)$ và $\vec{b}=(1;0;3)$. Tọa độ của vectơ tích có hướng $[\vec{a},\vec{b}]$ là:}{%
    \begin{tasks}(4)
        \task $(3;-5;-1)$
        \task $(3;5;-2)$
        \task $(2;-3;-1)$
        \task $(2;3;-1)$
    \end{tasks}
}

\questionLinesManual{8}{77}{Cho hai vectơ $\vec{a}=(2;2;-4)$, $\vec{b}=(1;1;-2)$. Mệnh đề nào sau đây sai?}{%
    \begin{tasks}(4)
        \task $[\vec{a},\vec{b}]=\vec{0}$
        \task $[\vec{a},\vec{b}]\ne\vec{0}$
        \task $|\vec{a}|=2|\vec{b}|$
        \task $\vec{a}=2\vec{b}$
    \end{tasks}
}

\questionLinesManual{8}{78}{Cho ba điểm $A(1;1;0)$, $B(0;-1;1)$, $C(1;2;1)$. Diện tích tam giác $ABC$ bằng:}{%
    \begin{tasks}(4)
        \task $\sqrt{11}$
        \task $\dfrac{1}{2}$
        \task $\dfrac{\sqrt{11}}{2}$
        \task $\dfrac{3}{2}$
    \end{tasks}
}

\questionLinesManual{8}{79}{Cho hình bình hành $ABCD$ có $A(1;0;1)$, $B(2;1;2)$ và giao điểm hai đường chéo $I\left(\dfrac{3}{2}; 0; \dfrac{3}{2}\right)$. Diện tích của hình bình hành $ABCD$ bằng:}{%
    \begin{tasks}(4)
        \task $\sqrt{2}$
        \task $\sqrt{5}$
        \task $\sqrt{6}$
        \task $\sqrt{3}$
    \end{tasks}
}

\questionLinesManual{8}{80}{Cho hình hộp $ABCD.A'B'C'D'$ có $\vec{AB}=(2;0;0)$, $\vec{AD}=(0;4;0)$ và $\vec{AA'}=(0;1;3)$. Tọa độ của vectơ $\vec{AC'}$ là:}{%
    \begin{tasks}(4)
        \task $(2;4;3)$
        \task $(2;3;5)$
        \task $(2;5;3)$
        \task $(1;5;3)$
    \end{tasks}
}

\questionLinesManual{8}{81}{Cho các điểm $A(1;0;0)$, $B(0;1;0)$, $C(0;0;1)$ và $D(-2;1;-1)$. Thể tích của khối tứ diện $ABCD$ bằng:}{%
    \begin{tasks}(4)
        \task $2$
        \task $1$
        \task $\dfrac{1}{3}$
        \task $\dfrac{1}{2}$
    \end{tasks}
}

\questionLinesManual{8}{82}{Cho tứ diện $ABCD$ có $A(-1;2;1)$, $B(0;0;-2)$, $C(1;0;1)$, $D(2;1;-1)$. Thể tích tứ diện $ABCD$ bằng:}{%
    \begin{tasks}(4)
        \task $\dfrac{1}{3}$
        \task $\dfrac{2}{3}$
        \task $\dfrac{4}{3}$
        \task $\dfrac{8}{3}$
    \end{tasks}
}

\questionLinesManual{8}{83}{Cho hình hộp $ABCD.A'B'C'D'$ có $A(1;1;-6)$, $B(0;0;-2)$, $C(-5;1;2)$ và $D'(2;1;-1)$. Thể tích khối hộp đã cho bằng:}{%
    \begin{tasks}(4)
        \task $12$
        \task $19$
        \task $38$
        \task $42$
    \end{tasks}
}

\questionLinesManual{8}{84}{Cho hai điểm $A(1;-2;3)$ và $B(2;-3;4)$. Tìm điểm $M \in (Oxy)$ sao cho ba điểm $A, B, M$ thẳng hàng:}{%
    \begin{tasks}(4)
        \task $M(1;1;0)$
        \task $M(3;-4;5)$
        \task $M(-3;5;0)$
        \task $M(-2;1;0)$
    \end{tasks}
}

\questionLinesManual{8}{85}{Cho hai điểm $A(2;-2;1)$ và $B(0;1;2)$. Tọa độ điểm $M$ thuộc mặt phẳng $(Oxy)$ sao cho ba điểm $A, B, M$ thẳng hàng là:}{%
    \begin{tasks}(4)
        \task $M(2;-3;0)$
        \task $M(4;-5;0)$
        \task $M(4;5;0)$
        \task $M(0;0;1)$
    \end{tasks}
}

\questionLinesManual{8}{86}{Cho ba điểm $A(2;-1;5)$, $B(5;-5;7)$ và $M(x;y;1)$. Với giá trị nào của $x$ và $y$ thì ba điểm $A, B, M$ thẳng hàng?}{%
    \begin{tasks}(4)
        \task $x=4;\ y=7$
        \task $x=-4;\ y=-7$
        \task $x=4;\ y=-7$
        \task $x=-4;\ y=7$
    \end{tasks}
}

\questionLinesManual{8}{87}{Cho ba điểm $A(3;-4;0)$, $B(0;2;4)$, $C(4;2;1)$. Tìm tọa độ điểm $D$ thuộc trục $Ox$ sao cho $AD=BC$:}{%
    \begin{tasks}(2)
        \task $D(0;0;0)$ hoặc $D(6;0;0)$
        \task $D(0;-6;0)$
        \task $D(0;0;0)$ hoặc $D(-6;0;0)$
        \task $D(6;0;0)$
    \end{tasks}
}

\questionLinesManual{8}{88}{Cho ba điểm $A(3;2;1)$, $B(1;-1;2)$, $C(1;2;-1)$. Tìm tọa độ điểm $M$ thỏa mãn $\vec{OM}=2\vec{AB}-\vec{AC}$:}{%
    \begin{tasks}(4)
        \task $M(5;5;0)$
        \task $M(-2;-6;4)$
        \task $M(-2;6;-4)$
        \task $M(2;-6;4)$
    \end{tasks}
}

\questionLinesManual{8}{89}{Cho ba điểm $A(-3;1;-2)$, $B(-1;-1;-1)$, $C(-3;1;1)$. Độ dài của vectơ $\vec{AB}+2\vec{AC}$ bằng:}{%
    \begin{tasks}(4)
        \task $\sqrt{57}$
        \task $\sqrt{7}$
        \task $3\sqrt{33}$
        \task $\sqrt{17}$
    \end{tasks}
}

\questionLinesManual{8}{90}{Cho các véctơ $\vec{a}=(-2;0;3)$, $\vec{b}=(0;4;-1)$ và $\vec{c}=(m-2;m^2;5)$. Tìm $m$ để ba véctơ $\vec{a}, \vec{b}, \vec{c}$ đồng phẳng:}{%
    \begin{tasks}(2)
        \task $m=2$ hoặc $m=-4$
        \task $m=-2$ hoặc $m=-4$
        \task $m=-2$ hoặc $m=4$
        \task $m=1$ hoặc $m=6$
    \end{tasks}
}

\questionLinesManual{8}{91}{Cho bốn điểm $O(0;0;0)$, $A(0;1;-2)$, $B(1;2;1)$, $C(4;3;m)$. Tìm $m$ để bốn điểm $O, A, B, C$ đồng phẳng:}{%
    \begin{tasks}(4)
        \task $m=-7$
        \task $m=-14$
        \task $m=14$
        \task $m=7$
    \end{tasks}
}

\questionLinesManual{8}{92}{Cho tứ diện $ABCD$ biết $A(0;-1;3)$, $B(2;1;0)$, $C(-1;3;3)$, $D(1;-1;-1)$. Chiều cao $AH$ của tứ diện $ABCD$ bằng:}{%
    \begin{tasks}(4)
        \task $AH=\dfrac{\sqrt{29}}{2}$
        \task $AH=\dfrac{14}{\sqrt{29}}$
        \task $AH=\sqrt{29}$
        \task $AH=\dfrac{1}{\sqrt{29}}$
    \end{tasks}
}

\questionLinesManual{8}{93}{Cho hình hộp $ABCD.A'B'C'D'$. Biết $A(1;0;1)$, $B'(2;1;2)$, $D'(1;-1;1)$, $C(4;5;-5)$. Đỉnh $A'$ có tọa độ $(a;b;c)$. Giá trị của $2a+b+c$ bằng:}{%
    \begin{tasks}(4)
        \task $3$
        \task $7$
        \task $2$
        \task $8$
    \end{tasks}
}

\questionLinesManual{8}{94}{Cho tam giác $ABC$ có $A(1;2;-1)$, $B(2;-1;3)$, $C(-4;7;5)$. Tọa độ chân đường phân giác trong góc $B$ của tam giác $ABC$ là:}{%
    \begin{tasks}(4)
        \task $\left(\dfrac{11}{3}; -2; 1\right)$
        \task $(-2; 11; 1)$
        \task $\left(\dfrac{2}{3}; \dfrac{11}{3}; \dfrac{1}{3}\right)$
        \task $\left(-\dfrac{2}{3}; \dfrac{11}{3}; 1\right)$
    \end{tasks}
}

\questionLinesManual{8}{95}{Cho tam giác $ABC$ với $A(1;2;-1)$, $B(2;-1;3)$, $C(-4;7;5)$. Độ dài đường phân giác trong kẻ từ đỉnh $B$ của tam giác $ABC$ là:}{%
    \begin{tasks}(4)
        \task $\dfrac{2\sqrt{74}}{5}$
        \task $\dfrac{2\sqrt{74}}{3}$
        \task $\dfrac{3\sqrt{73}}{3}$
        \task $2\sqrt{30}$
    \end{tasks}
}

\questionLinesManual{8}{96}{Trong không gian $Oxyz$, cho hai điểm $A(1;1;0)$, $B(2;-1;2)$. Tọa độ điểm $M$ thuộc trục $Oz$ sao cho $MA^2+MB^2$ đạt giá trị nhỏ nhất là:}{%
    \begin{tasks}(4)
        \task $M(0;0;-1)$
        \task $M(0;0;0)$
        \task $M(0;0;2)$
        \task $M(0;0;1)$
    \end{tasks}
}

\questionLinesManual{8}{97}{Cho tứ diện $ABCD$ có $A(2;1;-1)$, $B(3;0;1)$, $C(2;-1;3)$ và điểm $D$ nằm trên trục $Oy$. Biết thể tích tứ diện $ABCD$ bằng $5$. Tọa độ điểm $D$ là:}{%
    \begin{tasks}(2)
        \task $D(0;-7;0)$
        \task $D(0;8;0)$
        \task $D(0;7;0)$ hoặc $D(0;-8;0)$
        \task $D(0;-7;0)$ hoặc $D(0;8;0)$
    \end{tasks}
}

\questionLinesManual{8}{98}{Cho bốn điểm $A(2;5;1)$, $B(-2;-6;2)$, $C(1;2;-1)$ và điểm $D(d;d;d)$. Tìm tham số $d$ để biểu thức $|\vec{DB}-2\vec{AC}|$ đạt giá trị nhỏ nhất:}{%
    \begin{tasks}(4)
        \task $d=3$
        \task $d=4$
        \task $d=1$
        \task $d=2$
    \end{tasks}
}

\questionLinesManual{8}{99}{Cho ba điểm $A(2;5;1)$, $B(-2;-6;2)$, $C(1;2;-1)$. Điểm $M$ thay đổi thỏa mãn biểu thức $MA^2-MB^2-MC^2$ đạt giá trị lớn nhất. Khi đó độ dài đoạn thẳng $OM$ bằng:}{%
    \begin{tasks}(4)
        \task $3\sqrt{10}$
        \task $3\sqrt{5}$
        \task $3\sqrt{3}$
        \task $2\sqrt{3}$
    \end{tasks}
}

\questionLinesManual{8}{100}{Cho hình hộp chữ nhật $ABCD.A'B'C'D'$ có $A$ trùng với gốc tọa độ $O$, các đỉnh $B(m;0;0)$, $D(0;m;0)$, $A'(0;0;n)$ với $m,n>0$ và $m+n=4$. Gọi $M$ là trung điểm của cạnh $CC'$. Thể tích tứ diện $BDAM$ đạt giá trị lớn nhất bằng:}{%
    \begin{tasks}(4)
        \task $\dfrac{245}{108}$
        \task $\dfrac{9}{4}$
        \task $\dfrac{64}{27}$
        \task $\dfrac{75}{32}$
    \end{tasks}
}

\vfill
\begin{center}
    \textbf{--- HẾT BÀI TẬP VỀ NHÀ (100 CÂU) ---}
\end{center}

\end{document}